\documentclass[11pt,letterpaper,twoside]{article}
\usepackage[margin=1in,headheight=40pt]{geometry}
\usepackage{andyspackage}
\usepackage[header=book,font=times,accent=purple]{andysnotes}
\title{Box and Page Break Demonstration}
\author{Andrew M. Hagy}
\date{}
\begin{document}
\setcounter{page}{0}
\AndyMakeTitle
\section{Long statements}
\subsection{Shared numbering}
\begin{definition}A brief unboxed definition.\end{definition}
\begin{AsideBlock}An aside in the same sequence.\end{AsideBlock}
\begin{TheoremBlock}[A long theorem with a nested proof]
A theorem can continue across a page. The following text intentionally
fills enough space to test the box at a page boundary.

\foreach \n in {1,...,13}{%
This paragraph explains that every statement remains visible while the
page builder is free to choose a break between lines. The mathematical
content here is only a layout test. We can use $\Gamma\entails\varphi$,
$\NN$, $\RR$, and $A\containedin B$ without adding another package.\par}

\begin{ProofBlock}
\foreach \n in {1,...,22}{%
The proof is nested in the theorem. This line tests a page break within
or near the proof while preserving a visible start and finish.\par}
\end{ProofBlock}
\end{TheoremBlock}
\begin{CorollaryBlock}A short corollary follows its theorem.\end{CorollaryBlock}
\end{document}
